\begin{problem}[Generic points through quotient spectra]
Let $I\subseteq A$ and let
\[
\pi:\Spec(A/I)\xrightarrow{\sim}V_A(I)
\]
be the quotient-spectrum homeomorphism.  Prove that $\pi$ preserves specialization,
closed points relative to $V(I)$, irreducible closed subsets, and their generic points.
\end{problem}

\paragraph{Why this problem matters.}
It shows that passing to a closed subspace does not change the internal pointwise geometry.

\paragraph{Useful ideas before starting.}
Primes of $A/I$ are exactly $\mathfrak p/I$ with $\mathfrak p\supseteq I$, and inclusion
is preserved.

\paragraph{Guided hints.}
Write
\[
\mathfrak p/I\subseteq\mathfrak q/I
\iff
\mathfrak p\subseteq\mathfrak q.
\]

\begin{solution}
The map sends $\mathfrak p/I$ to $\mathfrak p$.  Since
\[
\mathfrak p/I\subseteq\mathfrak q/I
\iff
\mathfrak p\subseteq\mathfrak q,
\]
it preserves the specialization order.  Because it is a homeomorphism onto $V(I)$, it
also preserves closed subsets and irreducibility.

If $Z'\subseteq\Spec(A/I)$ is irreducible closed with generic point $\mathfrak p/I$, then
\[
\overline{\{\mathfrak p/I\}}=Z'.
\]
Applying $\pi$ gives
\[
\overline{\{\mathfrak p\}}^{\,V(I)}=\pi(Z'),
\]
so $\mathfrak p$ is the generic point of the corresponding irreducible closed subset of
$V(I)$.  The converse follows by $\pi^{-1}$.

A point is closed relative to $V(I)$ exactly when its preimage is closed in
$\Spec(A/I)$.
\end{solution}

\paragraph{What happened mathematically.}
The quotient correspondence is not merely a bijection of primes; it transports the entire
generic/closed/specialization structure.

\paragraph{Extensions.}
Express the generic points of $V(I)$ directly in terms of minimal primes over $I$.

\paragraph{Provenance.}
New synthesis problem connecting VI/07 quotient spectra with the present pointwise theory.
